\documentclass[preprint,12pt,authoryear]{elsarticle}

\usepackage{amsmath,amssymb,amsfonts}
\usepackage{amsthm}
\let\origAND\AND
\let\AND\relax
\usepackage{algorithmic}
\usepackage{graphicx}
\usepackage{textcomp}
\usepackage{xcolor}
\usepackage{booktabs}
\usepackage{geometry}
\geometry{a4paper, margin=1in}
\usepackage{microtype}
\emergencystretch=3em

\usepackage{hyperref}
\hypersetup{
    colorlinks=true,
    linkcolor=blue,
    citecolor=blue,
    filecolor=blue,
    urlcolor=blue
}

\graphicspath{{figures/}{./}}

\journal{Engineering Applications of Artificial Intelligence}

\newtheorem{theorem}{Theorem}
\newtheorem{lemma}{Lemma}
\newtheorem{definition}{Definition}

\begin{document}

\begin{frontmatter}

\title{An Autonomous Multi-Agent Spatial Decision Support System for EV Charging Infrastructure Siting via Stochastic $M/M/c$ Queueing and Heterogeneous Data Fusion}

\author[inst1]{Harsh Pradeepkumar Ambule\corref{cor1}}
\ead{harshambule612@gmail.com}
\cortext[cor1]{Corresponding author}

\affiliation[inst1]{organization={Department of Artificial Intelligence, Priyadarshini Bhagwati College of Engineering (PBCOE)},
            city={Nagpur},
            state={Maharashtra},
            postcode={440024},
            country={India}}

\begin{abstract}
Rapid global adoption of electric vehicles (EVs) necessitates scalable computational frameworks for charging network operators (CPOs) and urban planners to identify optimal charging hub locations and forecast driver wait times. Existing infrastructure decision-support tools suffer from three fundamental limitations: (1) single-source geospatial data silos that miss live port hardware telemetry, (2) deterministic assumptions that underestimate stochastic queue wait times, and (3) black-box generative AI wrappers prone to spatial hallucination. In this paper, we present \textbf{EV-Spatial-Intelligence-Engine}, an autonomous multi-agent spatial decision-support system. The platform orchestrates six specialized asynchronous agents: an \textit{Orchestrator}, \textit{DriverAssistant}, \textit{Advisor}, \textit{Data}, \textit{Scoring}, and an \textit{ExplanationAgent} powered by Vertex AI Gemini 2.0 Flash. The engine integrates an analytical multi-server $M/M/c$ Erlang C queueing formulation to estimate median ($p_{50}$) and 90th percentile ($p_{90}$) driver wait times, alongside a Gradient Boosted Tree regressor for port utilization forecasting. We prove and numerically verify a percentile-ordering invariant ($p_{50}\le p_{90}$) across 728 distinct arrival-rate/service-rate/capacity configurations with zero violations. Against a 25{,}000-session Monte Carlo discrete-event simulation ground truth, the proposed $M/M/c$ model achieves a tail wait-time Mean Absolute Error (MAE) of 9.39 minutes, outperforming a lumped $M/M/1$ approximation (13.67 min MAE, +45.6\%) and a static duration heuristic (38.45 min MAE, +309.4\%). A 3-tier data provenance framework tags every spatial feature with an explicit confidence label (\texttt{live}, \texttt{estimated}, \texttt{fallback}) to ground the explanation layer and prevent unlabeled or fabricated attributes from reaching the end user. The engine is deployed as a containerized service on Google Cloud Run and validated with a 53-check automated harness spanning API reliability, input-validation robustness, cross-continental spatial coverage, geodesic accuracy, and the queueing invariant above; we additionally report a controlled synthetic sensitivity analysis illustrating how a multi-attribute siting utility differs from density-only heuristics under varying grid-capacity and competition conditions. We discuss the scope and limitations of the current evaluation and outline the field validation required before the siting-quality and hallucination-mitigation claims can be made empirically.
\end{abstract}

\begin{keyword}
Multi-Agent Systems \sep Electric Vehicle Infrastructure \sep $M/M/c$ Queueing Theory \sep Spatial Decision Support \sep Explainable Artificial Intelligence \sep Heterogeneous Data Fusion
\end{keyword}

\end{frontmatter}

\section{Introduction}
\label{sec:intro}
The decarbonization of terrestrial transportation relies fundamentally on the rapid expansion of public electric vehicle (EV) charging networks. However, charging network operators (CPOs) and municipal transportation authorities face substantial financial and operational risks when deploying high-power Direct Current Fast Charging (DCFC) hubs. Siting decisions require balancing competing multi-dimensional constraints: localized commercial demand, high-voltage electrical grid capacity, land acquisition capital expenditures, surrounding amenity density, and competitive market saturation.

Current industry approaches and academic prototypes typically exhibit three systemic failure modes:
\begin{enumerate}
    \item \textbf{The Single-Source Data Silo:} Systems relying solely on isolated proprietary APIs or single open databases operate with severe spatial blind spots, missing live socket states and utility constraints.
    \item \textbf{Deterministic Queueing Simplification:} Planners frequently assume zero wait times or constant average charging durations. In practice, arrivals follow stochastic Poisson processes, causing systematic underestimation of tail wait-time delays ($p_{90}$) during peak traffic.
    \item \textbf{Generative AI Hallucinations:} Recent efforts to leverage Large Language Models (LLMs) for infrastructure explanation often employ ungrounded zero-shot prompting, which is known in the broader LLM literature to fabricate unsupported details when inputs are incomplete \citep{lewis2020retrieval}.
\end{enumerate}

To resolve these challenges, this paper presents \textbf{EV-Spatial-Intelligence-Engine}, a decoupled, 4-layer autonomous multi-agent decision support architecture. The core contributions of this work are:
\begin{itemize}
    \item A decoupled 6-agent asynchronous coordination pipeline for spatial discovery, multi-criteria scoring, and natural-language explanation.
    \item A closed-form $M/M/c$ Erlang C percentile formulation, with a proven and numerically verified $p_{50}\le p_{90}$ ordering invariant across 728 parameter configurations.
    \item A 3-tier data provenance labeling scheme (\texttt{live} / \texttt{estimated} / \texttt{fallback}) designed to make every downstream inference auditable, including by the natural-language explanation layer.
    \item An automated, repository-reproducible evaluation harness (53 checks) covering API reliability, input validation, cross-continental spatial coverage, geodesic accuracy, and the queueing invariant, together with a discrete-event-simulation comparison of queueing accuracy against two standard baselines.
\end{itemize}
We deliberately scope the empirical claims of this paper to what the accompanying artifact reproduces end-to-end (Section~\ref{sec:limitations}); siting-quality and hallucination-mitigation claims are presented as a design rationale supported by a controlled synthetic analysis, not as field-validated results, pending the study described in Section~\ref{sec:limitations}.

\section{Related Work}
\label{sec:related}
Classical EV charging station site-selection has historically been formulated as mathematical programming problems, such as the Maximal Covering Location Problem (MCLP) or the Flow-Refueling Location Model \citep{kempton2005electric}. Multi-Agent Systems (MAS) provide a natural computational paradigm for decentralized cyber-physical coordination in energy systems \citep{morstyn2020multiagent}, with recent advances applying multi-agent deep reinforcement learning to EV charging station placement and capacity planning \citep{zhang2023multiagent}. Furthermore, spatial multi-criteria decision-making frameworks \citep{chen2024spatial} incorporate multi-attribute evaluation techniques such as the Analytic Hierarchy Process \citep{saaty1988analytic} and Pareto multi-objective evolutionary algorithms \citep{deb2002fast} to balance conflicting urban planning metrics. In parallel, short-term demand forecasting using spatial-temporal gradient boosting trees \citep{wang2023gradient} and big data technologies \citep{arias2016forecasting} has demonstrated high predictive fidelity across heterogeneous urban regions. Open geospatial data platforms such as OpenStreetMap \citep{haklay2008openstreetmap} and programmatic street network tools like OSMnx \citep{boeing2017osmnx} have significantly lowered the barriers to scalable geospatial feature extraction.

Queueing theory forms the standard mathematical foundation for service-facility wait-time modeling \citep{erlang1917solution,shortle2018fundamentals}. Prior MAS implementations for this domain are typically evaluated only in simulation, without a containerized production deployment, and published evaluations most often report expected average wait times rather than tail percentiles, even though tail delay ($p_{90}$, $p_{99}$) is what drives driver dissatisfaction and queue instability. Pricing- and routing-based approaches to differentiated EV charging service have also been studied as an alternative or complementary lever to physical siting \citep{moradipari2020pricing}. Our work targets the tail-percentile gap by deriving and verifying exact percentile expressions across a wide parameter grid.

\section{System Architecture and Multi-Agent Orchestration}
\label{sec:architecture}
The proposed platform is organized into four operational layers, illustrated in Figure~\ref{fig:architecture}.
\begin{enumerate}
    \item \textbf{Interface \& Gateway Layer:} An asynchronous FastAPI gateway \citep{tiangolo2018fastapi} coupled with an interactive Leaflet.js single-page application, deployed on containerized Google Cloud Run.
    \item \textbf{Autonomous Multi-Agent Orchestration Layer:} Six specialized agents --- \textit{OrchestratorAgent} (request lifecycle and routing), \textit{DriverAssistantAgent} (spatial search and connector-taxonomy normalization), \textit{AdvisorAgent} (administrative-district resolution via Nominatim/Overpass), \textit{DataAgent} (BigQuery telemetry and OCPP 1.6 WebSocket hardware sync), \textit{ScoringAgent} (multi-attribute utility scoring, weights validated at start-up to sum to $1.0\pm0.01$), and \textit{ExplanationAgent} (Vertex AI Gemini 2.0 Flash \citep{team2024gemini}, built upon transformer attention foundations \citep{vaswani2017attention,achiam2023gpt4}, grounded strictly in provenance-tagged inputs).
    \item \textbf{Mathematical \& Predictive Layer:} The $M/M/c$ Erlang C queueing model (Section~\ref{sec:math}), a Gradient Boosted Tree demand regressor, and an automated invariant-verification mesh.
    \item \textbf{Live Data Fusion Layer:} OpenChargeMap, OpenStreetMap/Overpass \citep{haklay2008openstreetmap}, Google Places, NREL AFDC, BigQuery, and OCPP 1.6, merged via a Haversine-distance provider-deduplication engine.
\end{enumerate}

\begin{figure}[htbp]
\centering
\includegraphics[width=0.85\textwidth]{fig1_architecture.pdf}
\caption{Four-layer system architecture: interface/gateway, multi-agent orchestration, mathematical/predictive modeling, and live data fusion.}
\label{fig:architecture}
\end{figure}

\section{Mathematical Formulation}
\label{sec:math}
\subsection{Stochastic $M/M/c$ Erlang C Model}
Consider an EV fast-charging hub equipped with $c \in \mathbb{N}^+$ identical charging ports. Vehicle arrivals follow a Poisson process with arrival rate $\lambda$, and service times are exponentially distributed with service rate $\mu$ per port (mean service duration $1/\mu$).

The traffic intensity per port is defined as:
\begin{equation}
\rho = \frac{\lambda}{c \cdot \mu}, \quad (\rho < 1)
\end{equation}

The steady-state probability that an arriving vehicle finds all $c$ charging ports occupied and must queue is given by Erlang's C formula \citep{erlang1917solution}:
\begin{equation}
C(c, a) = \frac{\frac{a^c}{c!(1-\rho)}}{\sum_{k=0}^{c-1} \frac{a^k}{k!} + \frac{a^c}{c!(1-\rho)}}, \quad \text{where } a = \frac{\lambda}{\mu}
\end{equation}

The unconditional cumulative distribution function (CDF) of waiting time $W_q$ in the queue is:
\begin{equation}
P(W_q \le t) = 1 - C(c, a) \cdot e^{-c\mu(1-\rho)t}, \quad t \ge 0
\end{equation}

Solving for the $p$-th percentile wait time $t_p$ analytically yields:
\begin{equation}
t_p = \max\left(0, \; -\frac{\ln\left(\frac{1 - p}{C(c, a)}\right)}{c\mu(1-\rho)}\right)
\end{equation}

\noindent\textbf{Theorem 1 (Percentile Invariant Ordering).} \textit{For any stable queue configuration with $\rho < 1$, the median wait time $t_{0.50}$ and the 90th percentile wait time $t_{0.90}$ satisfy $t_{0.50} \le t_{0.90}$ unconditionally.}

\textit{Proof.} If $C(c, a) \le 0.10$, then $t_{0.50} = t_{0.90} = 0$. If $0.10 < C(c, a) \le 0.50$, then $t_{0.50} = 0 < t_{0.90}$. If $C(c, a) > 0.50$, $t_{0.50} = \frac{\ln(2C)}{c\mu(1-\rho)}$ and $t_{0.90} = \frac{\ln(10C)}{c\mu(1-\rho)}$. Since $\ln(10C) > \ln(2C)$ for all $C>0$, $t_{0.50} < t_{0.90}$. Thus $t_{0.50} \le t_{0.90}$ holds for all valid parameter spaces. $\square$

\textbf{Numerical verification.} We evaluate the invariant over a full factorial grid of $\lambda \in \{0.5,1,2,3,4,5,6,7,8,10,12,15,20\}$ (13 values), mean session length $\in\{10,15,20,25,30,40,60\}$ minutes (7 values), and $c\in\{1,2,3,4,5,6,8,10\}$ ports (8 values) --- $13\times7\times8 = 728$ configurations. We independently re-executed this sweep against the released \texttt{models/queue\_model.py} and confirmed \textbf{0 violations (100.0\% adherence)}, matching Section~\ref{sec:experiments}.

\begin{figure}[htbp]
\centering
\includegraphics[width=0.7\textwidth]{fig2_queue_flow.pdf}
\caption{Data flow through the $M/M/c$ queue model, from raw arrival/service parameters to the provenance-tagged $p_{50}$/$p_{90}$ output consumed by the explanation layer. Source: \texttt{assets/queue-model-flow.svg}, restyled per the note in Figure~\ref{fig:architecture}.}
\label{fig:queueflow}
\end{figure}

\section{Live Data Fusion \& 3-Tier Provenance}
\label{sec:provenance}
Every station and zone recommendation carries an immutable provenance label (Figure~\ref{fig:provenance}): \textbf{LIVE} (confidence 0.95--1.00, from active OCPP 1.6 telemetry or verified BigQuery port sessions), \textbf{ESTIMATED} (0.70--0.94, from multi-source API fusion, the $M/M/c$ model, or the GBT forecaster), and \textbf{FALLBACK} (0.30--0.69, deterministic seeded estimators used when upstream APIs fail, always accompanied by an explicit \texttt{fallback\_reason} string). This design is intended to prevent the \textit{ExplanationAgent} from presenting an unverified or missing value as fact: the LLM is prompted to condition its language on the tier of the data it is describing, so a \texttt{fallback}-tier attribute is described with appropriate hedging rather than stated as measured fact. We regard rigorous quantification of how much this reduces hallucinated output, relative to an ungrounded baseline, as an open empirical question requiring a real LLM-in-the-loop study (Section~\ref{sec:limitations}) rather than a claim this paper currently supports with data.

\begin{figure}[htbp]
\centering
\includegraphics[width=0.9\textwidth]{fig3_data_provenance.pdf}
\caption{The 3-tier data provenance framework: source, confidence range, and the explicit \texttt{fallback\_reason} contract for Tier 3. Source: \texttt{assets/data-provenance-matrix.svg}, restyled per Figure~\ref{fig:architecture}.}
\label{fig:provenance}
\end{figure}

\section{Experimental Results}
\label{sec:experiments}

\subsection{Queue Wait-Time Error vs. Discrete-Event Simulation Ground Truth}
We evaluate wait-time predictions across five traffic-intensity regimes ($\rho \in [0.25, 0.90]$) for a 4-port fast-charging hub ($c=4$, $\mu=2.0$ sessions/hr) against a 25{,}000-session Monte Carlo discrete-event simulation (DES) ground truth, using the earliest-available-port discipline. Table~\ref{tab:queue_comparison} and Figure~\ref{fig:queuecomparison} report the result; both were regenerated independently from \texttt{eval/comparative\_baseline.py} for this revision (seed = 42) and reproduce the original figures exactly.

\begin{table}[htbp]
\centering
\caption{Comparison of 90th-percentile ($p_{90}$) wait-time estimation errors across traffic intensities. Reproducible via \texttt{python eval/comparative\_baseline.py}.}
\label{tab:queue_comparison}
\resizebox{\linewidth}{!}{%
\begin{tabular}{cccccc}
\toprule
$\lambda$ (veh/h) & Intensity $\rho$ & DES Ground Truth & Proposed $M/M/c$ & Lumped $M/M/1$ & Static Heuristic \\
\midrule
2.0 & 0.25 & 0.00 min & \textbf{0.00 min} & 5.76 min & 15.00 min \\
4.0 & 0.50 & 7.84 min & \textbf{8.30 min} & 17.27 min & 15.00 min \\
5.5 & 0.69 & 31.56 min & \textbf{33.83 min} & 37.99 min & 15.00 min \\
6.5 & 0.81 & 62.05 min & \textbf{72.92 min} & 74.83 min & 15.00 min \\
7.2 & 0.90 & 121.47 min & \textbf{154.80 min} & 155.42 min & 15.00 min \\
\midrule
\textbf{MAE} & --- & --- & \textbf{9.39 min} & 13.67 min (+45.6\%) & 38.45 min (+309.4\%) \\
\textbf{RMSE} & --- & --- & \textbf{15.72 min} & 17.21 min & 53.10 min \\
\bottomrule
\end{tabular}%
}
\end{table}

\begin{figure}[htbp]
\centering
\includegraphics[width=0.75\textwidth]{fig4_queue_comparison.pdf}
\caption{$p_{90}$ wait time versus traffic intensity $\rho$: DES ground truth, the proposed $M/M/c$ model, a lumped $M/M/1$ baseline, and a static 15-minute heuristic. Generated directly from \texttt{eval/comparative\_baseline.py} output.}
\label{fig:queuecomparison}
\end{figure}

\textbf{Discussion.} The static heuristic fails to capture the non-linear tail-queue explosion under heavy traffic ($\rho \ge 0.80$), underestimating delay by over 100 minutes at $\rho=0.90$. The lumped $M/M/1$ model tracks the qualitative escalation but consistently overestimates wait time at moderate load due to server-pooling distortion. The proposed $M/M/c$ formulation tracks the DES ground truth closely across all tested regimes, with error concentrated at the highest-intensity point ($\rho=0.90$), where all analytical approximations of Erlang-C systems are known to be most sensitive to the exponential-service-time assumption.


\subsection{Offline Demand Forecaster Temporal Evaluation}
\label{sec:forecaster_eval}
To evaluate the Gradient Boosted Tree demand regressor within the predictive layer, we backtested the model using an offline chronological holdout split implemented in \texttt{models/forecaster\_eval.py}. To eliminate temporal data leakage inherent in random shuffling, the first 80\% of chronological sessions serve as the training set, while the remaining 20\% form the unseen test set. The model ingests cyclic sinusoidal/cosinusoidal encodings for hour-of-day, day-of-week, and month, alongside weekend flags, charger power ($\text{kW}$), and lag-1h utilization. Compared against an empirical hour-of-day mean baseline, the GBT forecaster reduces Mean Absolute Error from 0.0899 to 0.0012 (a 98.7\% relative error reduction) and RMSE from 0.1106 to 0.0017, confirming that dynamic feature conditioning yields substantial predictive fidelity over static historical lookup tables before parameters enter the /M/c$ queue solver.

\subsection{Automated Evaluation Harness}
\begin{figure}[htbp]
\centering
\includegraphics[width=0.75\textwidth]{fig6_global_latency.pdf}
\caption{Global metropolitan generalization and sub-second median discovery response latency ({50}$) across ten cities spanning five continents on Google Cloud Run.}
\label{fig:globallatency}
\end{figure}

The system ships an automated harness (\texttt{eval/benchmark.py}) designed to evaluate deployment health across 10 functional modules: API reliability (11 endpoints), input-validation correctness (10 malformed-payload cases), cross-continental spatial coverage (10 cities, 5 continents), data-quality labeling completeness, output type-filtering correctness, GeoJSON schema compliance (5 checks), geodesic (Haversine) distance accuracy against 5 known city-pair distances, the 728-configuration queueing invariant (5 boundary checks), multi-source data-fusion presence, and expansion-ranking monotonicity --- \textbf{53 checks in total}. In offline verification, all 10 mathematical and geodesic invariant checks pass with 100\% adherence (0 invariant violations across 728 queue regimes and Haversine distance errors $<0.05\%$), while the remaining 43 checks validate end-to-end HTTP responses against the live Cloud Run microservices. To guarantee permanent evaluation reproducibility independent of live cloud runtime hosting or external credentials, all 53 checks can also be executed locally offline by deploying the containerized microservices via \texttt{docker compose up} and running \texttt{python eval/benchmark.py} against \texttt{http://localhost:8000}.

\subsection{Synthetic Sensitivity Analysis: Multi-Attribute vs. Density-Only Siting}
\label{sec:siting}
To illustrate \emph{qualitatively} how a multi-attribute siting criterion differs from a density-only heuristic, we constructed 50 synthetic candidate parcels with independently sampled POI density, grid headroom, competing-charger count, and installation cost (uniform/Poisson distributions, seed = 101), and compared the top-5 parcels selected by (a) POI density alone and (b) a linear utility combining all four normalized attributes (weights 0.35/0.30/$-$0.20/$-$0.15). Figure~\ref{fig:siting} and Table~\ref{tab:siting_comparison} report the result.

\begin{table}[htbp]
\centering
\caption{Synthetic sensitivity analysis (not a field study): comparative performance on spatial siting viability metrics over synthetic candidate parcels.}
\label{tab:siting_comparison}
\resizebox{\linewidth}{!}{%
\begin{tabular}{lccc}
\toprule
Metric & Greedy POI Baseline & Multi-Attribute Utility & Difference \\
\midrule
Average Grid Headroom & 350.1 kW & \textbf{410.5 kW} & +17.3\% \\
Competitor Overlap & 3.4 stations & \textbf{2.2 stations} & $-$35.3\% \\
Composite Viability Score & 0.339 & \textbf{0.394} & +16.1\% \\
Average Installation CapEx & \$96,435 & \textbf{\$87,393} & $-$9.4\% \\
\bottomrule
\end{tabular}%
}
\end{table}

\begin{figure}[htbp]
\centering
\includegraphics[width=\textwidth]{fig5_siting_tradeoffs.pdf}
\caption{Synthetic sensitivity analysis, plotted as four independent panels (grid headroom in kW, competitor count, composite score, and installation CapEx in USD) rather than one shared axis, since the four quantities are not dimensionally comparable. Because the multi-attribute formula is defined to weight grid headroom, competitor count, and cost, while the density-only baseline is defined to ignore all three, this figure demonstrates that the two selection rules behave differently under synthetic data with those axes of variation --- it does \emph{not} demonstrate that the multi-attribute approach identifies better real-world sites, which requires field data (Section~\ref{sec:limitations}).}
\label{fig:siting}
\end{figure}

\textit{Important scope note:} this analysis uses a simplified 4-factor formula on synthetic data for illustration, and is \emph{not} the same computation as the deployed \texttt{ScoringAgent} (which uses 6 config-driven weighted factors over real zip-code-indexed zone, competition, and usage records). We present it here as a design-rationale illustration, not as evidence that the deployed system outperforms density-only siting in the field --- see Section~\ref{sec:limitations}.

\subsection{Component Ablation Study}
\label{sec:ablation}
To quantify the contribution of individual subsystems, we conducted a component ablation using \texttt{eval/ablation\_study.py} (released alongside this manuscript). Two of the four planned ablation axes are backed by numbers reproducible today; the other two require experiments beyond a code toggle and are reported as ongoing work rather than filled in with placeholder figures.

\begin{table}[htbp]
\centering
\caption{Component ablation: axes 1--2 are independently reproducible today via \texttt{eval/ablation\_study.py}; axes 3--4 require live-deployment and LLM-in-the-loop experiments not yet run (see notes below).}
\label{tab:ablation}
\resizebox{\linewidth}{!}{%
\begin{tabular}{lll}
\toprule
Ablated Component & Full System & Ablated \\
\midrule
Erlang-C queueing model $\to$ static heuristic & MAE 9.39 min & MAE 38.45 min (+309.4\%) \\
Multi-agent parallelism $\to$ sequential awaits & $\sim$205--210 ms$^\dagger$ & $\sim$2,600--3,000 ms$^\dagger$ (12--15$\times$ slower) \\
Multi-source fusion $\to$ single-source (OCM only) & \multicolumn{2}{l}{\textit{Requires a live run against the deployed API; see note (3) below.}} \\
3-tier data provenance $\to$ ungrounded LLM prompt & \multicolumn{2}{l}{\textit{Requires a real LLM-in-the-loop experiment; see note (4) below.}} \\
\bottomrule
\end{tabular}%
}
\end{table}

\noindent$^\dagger$\textit{The parallelism axis reproduces the real \texttt{asyncio.gather} concurrency structure in \texttt{agents/driver\_agent.py} (independent per-station enrichment coroutines, each awaiting the same sequence of I/O calls as the real \texttt{\_enrich\_one} method), with simulated I/O latency standing in for network calls this evaluation environment cannot reach. The qualitative result --- an order-of-magnitude latency reduction from concurrent execution --- is structural and repeatable; the exact millisecond values vary with the running machine's scheduler overhead (we observed 207~ms/2{,}646~ms in one environment and 205~ms/3{,}034~ms in another, both from the identical script and seed). Report a range, not a single decimal, if citing this axis.}

\textit{(3) Multi-source fusion:} requires an environment-variable toggle in the client fan-out (restrict to \texttt{ocm\_client} only) and two live runs of \texttt{eval/benchmark.py} against the deployed Cloud Run instance, with and without the toggle. Not run in this revision.

\textit{(4) Data provenance:} requires real, repeated Gemini calls under constructed missing-data conditions with a disclosed human- or rubric-based fabrication-scoring methodology --- this is a genuine experiment to design, not a configuration flag. Not run in this revision.

\section{Limitations and Required Follow-up Validation}
\label{sec:limitations}
We are explicit about what this paper's evaluation does and does not establish, since the manuscript accompanies a public code repository that any reviewer can execute:
\begin{itemize}
    \item \textbf{Queueing accuracy (Section~\ref{sec:experiments}.1) and the 728-configuration invariant} are independently reproducible from the released code and are the strongest empirical claims in the paper.
    \item \textbf{The siting comparison (Section~\ref{sec:siting})} is a synthetic illustration of a design principle, not a validation of the deployed \texttt{ScoringAgent} against real candidate-site data. A field validation would require real zone/competition/usage records for a defined metropolitan area, run through the actual deployed agent, compared against a documented CPO siting baseline.
    \item \textbf{LLM hallucination mitigation} is a design goal of the 3-tier provenance framework, motivated by the well-documented tendency of ungrounded LLMs to fabricate detail under incomplete context \citep{lewis2020retrieval}. We do not currently report a quantitative hallucination-rate reduction, because doing so credibly requires a real experiment: repeated LLM calls under systematically constructed missing-data conditions, with human- or rubric-scored fabrication detection. This is planned future work, not a result of the present paper.
    \item \textbf{The automated evaluation harness} distinguishes between deterministic mathematical invariants verified offline (100\% passing) and live cloud endpoint availability, which depends on runtime infrastructure hosting status.
\end{itemize}

\section{Conclusion}
\label{sec:conclusion}
This paper presented an autonomous multi-agent spatial decision-support system for EV charging infrastructure planning, combining a 6-agent asynchronous architecture, a stochastic $M/M/c$ Erlang C queueing formulation with a proven and numerically verified percentile-ordering invariant across 728 configurations, and a 3-tier data-provenance framework intended to ground natural-language explanation. Against Monte Carlo discrete-event simulation ground truth, the proposed queueing model reduces tail wait-time estimation error by 45.6\% relative to a lumped $M/M/1$ baseline. We have been explicit about which additional claims --- multi-criteria siting superiority and LLM hallucination mitigation --- remain design rationale supported by synthetic illustration rather than field-validated results, and we outline the follow-up studies needed to substantiate them empirically.


\section*{Author Biography}
\noindent\textbf{Harsh Pradeepkumar Ambule} is an Independent Researcher based in Nagpur, Maharashtra, India. He received the B.Tech degree in Artificial Intelligence from Priyadarshini Bhagwati College of Engineering (PBCOE), Nagpur, in 2026 (e-mail: \href{mailto:harshambule612@gmail.com}{harshambule612@gmail.com}). His research focuses on autonomous multi-agent systems, spatial decision-support architectures, queueing theory, and explainable AI in electric vehicle charging infrastructure. He is also the author of a multi-modal SAR-optical-topographic cross-attention framework for satellite cloud removal, currently under peer review, and co-author of two publications on a multi-agent, LLM-powered system for real-time fake news verification and credibility analysis, developed as his undergraduate final-year project.

\bibliographystyle{elsarticle-harv}
\bibliography{references}

\end{document}